% RESEARCH_DRAFT_SHA256: ed0b3c32bd4606c185ce5ae42362191e27e37b871a930dd7e83c2dee1815d259
% Layout and source-reference wrappers expanded for the web renderer.
\documentclass[11pt]{article}
\usepackage[margin=25mm]{geometry}
\usepackage{fontspec,unicode-math}
\setmainfont{Latin Modern Roman}
\setmathfont{Latin Modern Math}
\usepackage{amsmath,amssymb,mathtools,xurl,hyperref}
\hypersetup{colorlinks=true,urlcolor=blue,linkcolor=blue}
\setlength{\parindent}{0pt}
\setlength{\parskip}{5pt}
\setlength{\emergencystretch}{3em}
\newcommand{\sref}[2]{\hyperlink{source:#1:#2}{[S#2]}}
\newcommand{\cataloguescope}[1]{\paragraph{Source question.}#1}
\newcommand{\floor}[1]{\left\lfloor#1\right\rfloor}
\begin{document}

\section*{Repeated augmentation question: an exact star case and a failure of minimum-cut-only coverage}
\noindent UnsolvedMath ID 30004908; OWR-8415356-019. Research attempt dated 11 October 2026.
\paragraph{Result in scope.}This manuscript contains a source-specific attempt and complete proofs of the stated auxiliary calculations. It makes no new solution or novelty claim for the entire target.
\section{Original problem and preserved definition}
The following definition and source list are reproduced from the repository. Any earlier inspection dates and reported literature claims in this preserved text describe that source record, not additional verification by this attempt.
% BEGIN PRESERVED DEFINITION
\subsection{Definitions and mathematical statement}

\paragraph{Shared notation.}Standard number systems are $\mathbb N,\mathbb Z,\mathbb Q,\mathbb R,\mathbb C$. Specialized notation and index ranges follow the cited source and the definitions below.


In $\alpha$-multiplicative connectivity augmentation the input is a $k$-edge-connected graph and a set of candidate links with nonnegative costs. One selects minimum-cost links so that the augmented graph is at least $\lceil\alpha k\rceil$-edge-connected. Edge connectivity is the minimum number of edges in a nontrivial cut. An approximation factor below two means a polynomial-time algorithm with cost at most $\rho$ times optimum for some $\rho<2$.
\paragraph{Record scope and qualifications.}
The catalogue title is retained as metadata. Named certificates, contracts, bounds or frameworks in that title are not asserted as results of the original source extract.
The supplied registry flags this extraction as unrecoverable. It is displayed as a literal source fragment; the definitions provide context rather than a replacement formal statement.

Source-record limitation: Source flags the statement as unrecoverable: The official report confirms that this is a duplicate extraction of the multiplicative connectivity augmentation problem retained in OWR-8415356-018. No separate clean mathematical statement is supplied; this source ID preserves its provenance.

Related catalogue records: \href{https://mehmetmars7.github.io/Erdosproblems-llm-hunter/problem.html?type=open_problems&id=30004907}{30004907}. Separate source IDs are preserved for their provenance.
The source text below is rendered literally where necessary; damaged formulas and composite extracts are not silently repaired into a new mathematical claim.
The supplied catalogue marks this record as unrecoverable. Its extract and terminology are documented, but a complete original problem has not been reconstructed.
\paragraph{Statement recorded in the supplied source.}
Therefore, we propose first studying \textbackslash{}alpha-multiplicative connectivity augmentation for \textbackslash{}alpha \textbackslash{}leq 6/5 which may alternately be called polygon augmentation, following the naming convention of the problems tree augmentation and cactus augmentation in which the goal is to increase the connectivity of a graph by 1. The main open problem is to design a better-than-2 approximation for this problem for any constant \textbackslash{}alpha > 1. Open Problem: \textbackslash{}alpha-Multiplicative Connectivity Augmentation and Polygon Augmentation Nathan Klein \href{https://doi.org/10.4171/owr/2021/53}{[S2]}
\paragraph{Source status.}
The supplied catalogue labels this source record partially solved; that metadata was checked against the supplied snapshot on 2026-10-11. The label is preserved for provenance and is not an independent verification of present open status or mathematical correctness. The catalogue literature assessment is dated 2026-08-22.
\subsection{Short English statement}
The supplied fragment concerns better-Than-Two Approximation for Multiplicative Connectivity Augmentation. Its complete mathematical target is not recovered.
\subsection{Sources}
\begin{itemize}
\item[S1] Ulam AI and the UnsolvedMath Contributors, UnsolvedMath: A Curated Collection of Open Mathematics Problems, record 30004908 (OWR-8415356-019). Catalogue curation and metadata: CC BY 4.0. The source extract is retained; definitions, English summaries and qualifications are editorial additions.
\url{https://huggingface.co/datasets/ulamai/UnsolvedMath}.
\item[S2] Combinatorial Optimization (2022), Oberwolfach Reports, DOI 10.4171/owr/2021/53, pp. 2893--2954
\url{https://doi.org/10.4171/owr/2021/53}.
\item[S3] Weighted connectivity augmentation
\url{https://arxiv.org/abs/2209.07860}.
\item Source-specific Oberwolfach Reports licences, including CC BY-SA 4.0 where applicable, are retained; see THIRD\_PARTY\_NOTICES.md.
\end{itemize}
% END PRESERVED DEFINITION
\subsection{Verbatim supplied catalogue question}
Therefore, we propose first studying \textbackslash{}alpha-multiplicative connectivity augmentation for \textbackslash{}alpha \textbackslash{}leq 6/5 which may alternately be called polygon augmentation, following the naming convention of the problems tree augmentation and cactus augmentation in which the goal is to increase the connectivity of a graph by 1. The main open problem is to design a better-than-2 approximation for this problem for any constant \textbackslash{}alpha > 1. Open Problem: \textbackslash{}alpha-Multiplicative Connectivity Augmentation and Polygon Augmentation Nathan Klein

\section{Duplicate scope}
This fragment repeats the augmentation problem retained as 30004907.
The original paragraph on printed page 2945 supplies the base
connectivity, weighted candidate links and rounded target. We examine
a simple base class where the deficient cuts can be characterized
exactly, then test extending the rule of covering only the original
minimum cuts. This is a separate attempted route on the same
underlying question, not a new independent conjecture.

\paragraph{An exact unit-cost star augmentation.}
Let $G$ be a star with $h\ge2$ leaves. Its edge
connectivity is one. Suppose every pair of leaves is available
as a unit-cost link, and $1<\alpha\le2$, so the target
is $\lceil\alpha\rceil=2$. A selected set of links is
feasible exactly when each leaf is incident to at least one link.
To prove this, represent any cut by its side not containing the
center; that side is a nonempty leaf subset $S$. Its original
cut size is $|S|$. If $|S|\ge2$, the target is
already met. If $S$ is a singleton, it needs exactly one
crossing link, equivalent to that leaf's having positive link degree.
This proves necessity and sufficiency for every cut.

The optimum number of links is $\lceil h/2\rceil$.
Every link covers at most two leaves, giving this lower bound.
For even $h$, pair all leaves and take the resulting $h/2$
links. For odd $h\ge3$, pair $h-1$ leaves and connect
the remaining leaf to any paired leaf; this uses $(h+1)/2$
links and covers every leaf. The construction is polynomial and
exact, but uses the complete candidate-link set, unit costs and
the target two. It does not solve arbitrary weighted augmentation.

\paragraph{Original minimum cuts need not be the only deficient cuts.}
Keep a star with four leaves, but now take $\alpha=3$,
so the target is three. Select two parallel links between leaves
one and two, and two between leaves three and four. Each singleton
leaf cut increases from one to three, so all original minimum
cuts meet the target. The cut with side consisting of leaves
one and two still has only its two star edges: none of the
selected links crosses it. Its capacity is two, below target.
Thus covering every original minimum cut to its required new
size is not sufficient for general multiplicative augmentation.
All originally sub-target cuts matter, including those that were
not minimum cuts. The source allows candidate links as a multiset,
so the parallel links in this countercheck respect that model.

This example uses $\alpha=3$, outside the source's proposed
initial polygon range $\alpha\le6/5$. It tests the full
``any constant greater than one'' goal, not an impossibility claim
for that smaller range. Conversely the exact target-two star
construction lies within the smaller range when $\alpha\le6/5$
because the ceiling is still two. Neither observation transfers
its algorithm to arbitrary large base connectivity.

\section{Audit}
The primary augmentation paragraph and the duplicate link were
inspected. The audit checked the ceiling at base connectivity one,
both sides of every star cut, the odd-leaf construction and the
scope of the parallel-link example. All proofs are by explicit
cut counts, with no optimization solver used. The missing step is
an algorithm coordinating all deficient cuts with arbitrary costs
and base graph structure while achieving a uniform factor below
two. This remains unresolved under the repeated source record.

\section{Process and contributions}
Mathematics and model-assisted adversarial self-audit: OpenAI gpt-6-astra, ultra reasoning, as dispatched for this research batch; exact serving revision not exposed. Codex (GPT-6; exact serving revision not exposed) materially assisted with orchestration and format review. No human mathematical authorship, mathematical edits, or independent human proof checking is asserted.
Contribution prompt A was applied to the full definition and supplied source row, then prompt B was applied adversarially to the derived results. The item-specific audit above records the scope checks and rejected extrapolations. Local assembly and structural checks do not establish mathematical correctness. The source catalogue status is preserved in TOP\_PROBLEM and is separate from this attempt status.
\paragraph{Integration review.} Codex (GPT-6; exact serving revision not exposed) reviewed the written arguments and expanded layout and source-reference wrappers. Repository metadata and the production web renderer were checked using Python 3.13 and JavaScript with MathJax 3.2.2. This model-assisted review is not independent human proof verification; no PDF compilation or formal proof checking is asserted.

\section{Completion Estimate}
\noindent\textbf{COMPLETION: 0\%}
The proved material is a rederivation of known or auxiliary facts rather than a substantiated new advance on the original target. A complete-link unit-cost star case is solved exactly, but minimum-cut-only coverage fails in the general multiplicative target and no full approximation algorithm is supplied.
\section{References}
The complete original Sources subsection above is retained, including catalogue attribution and licence notices. The additional list records primary sources actually consulted for this attempt on 11 October 2026; uncited source links reproduced above were not all rechecked.
\begin{itemize}
\item Oberwolfach Report 53/2021, Combinatorial Optimization, original question and surrounding discussion read from the primary PDF text on 11 October 2026. Historical statements are not a current-status survey. \url{https://ems.press/content/serial-article-files/46931}.
\end{itemize}
\section{Final status}
\textbf{Unresolved in this attempt.} A complete-link unit-cost star case is solved exactly, but minimum-cut-only coverage fails in the general multiplicative target and no full approximation algorithm is supplied.
\end{document}
